% Original source preamble and source-level macro declarations.
% This archival file is not the blueprint driver preamble.

\documentclass[12 pt]{amsart}

\usepackage{hyperref}
\usepackage{etex}
\usepackage[shortlabels]{enumitem}
\usepackage{amsmath}
\usepackage{amsxtra}
\usepackage{amscd}
\usepackage{amsthm}
\usepackage{amsfonts}
\usepackage{amssymb}
\usepackage{eucal}
\usepackage[all]{xy}
\usepackage{graphicx}
\usepackage{tikz-cd}
\usepackage{mathrsfs}
\usepackage{subfiles}
\usepackage{mathpazo}
\usepackage[colorinlistoftodos, textsize=tiny]{todonotes}
\setlength{\marginparwidth}{2cm}
\usepackage{morefloats}
\usepackage{pdfpages}
\usepackage{thm-restate}
\usepackage[utf8]{inputenc}
\usepackage{epigraph}
\usepackage{csquotes}
\usepackage[margin=1.5in]{geometry}
\usepackage{adjustbox}
\usepackage{stmaryrd}
\usepackage{microtype}
\usepackage{scalerel}
\usepackage{stackengine}
\stackMath
\newcommand\reallywidehat[1]{%
\savestack{\tmpbox}{\stretchto{%
  \scaleto{%
    \scalerel*[\widthof{\ensuremath{#1}}]{\kern-.6pt\bigwedge\kern-.6pt}%
    {\rule[-\textheight/2]{1ex}{\textheight}}%WIDTH-LIMITED BIG WEDGE
  }{\textheight}% 
}{0.5ex}}%
\stackon[1pt]{#1}{\tmpbox}%
}
\parskip 1ex


%\usepackage{showlabels}

\graphicspath{ {images/} }

\RequirePackage{color}
\definecolor{myred}{rgb}{0.75,0,0}
\definecolor{mygreen}{rgb}{0,0.5,0}
\definecolor{myblue}{rgb}{0,0,0.65}

\usepackage{color}
\newcommand{\daniel}[1]{{\color{blue} \sf
    $\spadesuit\spadesuit\spadesuit$ DANIEL: [#1]}}
\newcommand{\rboxed}[1]{\boxed{#1}}
\newcommand{\aaron}[1]{{\color{red} \sf
    $\spadesuit\spadesuit\spadesuit$ AARON: [#1]}}



\usepackage{hyperref}
\hypersetup{citecolor=blue}
\usepackage{tikz}
\usetikzlibrary{matrix,arrows,decorations.pathmorphing}

\theoremstyle{plain}
\newtheorem{theorem}[subsubsection]{Theorem}
\newtheorem{proposition}[subsubsection]{Proposition}
\newtheorem{lemma}[subsubsection]{Lemma}
\newtheorem{corollary}[subsubsection]{Corollary}
\newtheorem{situation}[subsubsection]{Situation}
\newtheorem{problem}[subsubsection]{Problem}
\theoremstyle{definition}
\newtheorem{definition}[subsubsection]{Definition}
\newtheorem{remark}[subsubsection]{Remark}
\newtheorem{example}[subsubsection]{Example}
\newtheorem{exercise}[subsubsection]{Exercise}
\newtheorem{counterexample}[subsubsection]{Counterexample}
\newtheorem{convention}[subsubsection]{Convention}
\newtheorem{question}[subsubsection]{Question}
\newtheorem{conjecture}[subsubsection]{Conjecture} 
\newtheorem{goal}[subsubsection]{Goal}
\newtheorem{warn}[subsubsection]{Warning}
\newtheorem{fact}[subsubsection]{Fact}
\theoremstyle{remark}
\newtheorem{notation}[subsubsection]{Notation}
\newtheorem{construction}[subsubsection]{Construction}
\numberwithin{equation}{section}
\newcommand\nc{\newcommand}
\nc\on{\operatorname}
\nc\renc{\renewcommand}
\DeclareMathOperator\rk{rk}
\newcommand\se{\section}
\newcommand\ssec{\subsection}
\newcommand\sssec{\subsubsection}

\newcommand*{\sheafhom}{\mathscr{H}\kern -2pt om}
\newcommand*{\sheafend}{\mathscr{E}\kern -1pt nd}

\renewcommand{\sectionautorefname}{\S$\!$}
\renewcommand{\subsectionautorefname}{\S$\!$}
\renewcommand{\subsubsectionautorefname}{\S$\!$}


\title{Geometric local systems on very general curves and isomonodromy}
\author{Aaron Landesman, Daniel Litt}
\date{\today}


